\chapter*{Abstrak}
\addcontentsline{toc}{chapter}{Abstrak}

Ang pag-aaral ng kimika ay naglalaman ng mga pundamental na konsepto na kadalasang itinuturing na mahirap unawain ng mga mag-aaral, lalo na pagdating sa aplikasyon nito sa totoong buhay. Ang ganitong kahirapan ay maaaring magdulot ng cognitive overload, partikular kung kulang ang mga aralin sa praktikal at interaktibong gawain. Sa pamamagitan ng interaktibong pagkatuto, aktibong nahahasa ng mga mag-aaral ang kanilang kasanayan sa pagsusuri at mas nauunawaan nila kung paano nagagamit ang mga konsepto ng kimika sa pang-araw-araw na buhay.

Sinisiyasat ng pag-aaral na ito ang paggamit ng Game-Based Learning (GBL) sa pamamagitan ng pagbuo ng [GAME NAME HERE], isang pang-edukasyong seryosong laro na idinisenyo bilang pantulong na materyal sa pagkatuto para sa mga mag-aaral sa Baitang 6 sa ilalim ng MATATAG kurikulum sa Pilipinas. Nakatuon ang laro sa pisikal at kemikal na pagbabago ng materyal, na nagpapalawak sa kaalaman ng mga mag-aaral lampas sa tatlong karaniwang kaanyuan ng matter. Sa pamamagitan ng interaktibong gameplay, natututuhan ng mga mag-aaral ang mga konsepto tulad ng pagbabago ng estado at paglipat ng init, habang malinaw na natutukoy ang kaibahan ng pisikal at kemikal na pagbabago.

Layunin ng pag-aaral na ito na mapabuti ang pag-unawa, motibasyon, at emosyonal na pakikipag-ugnayan ng mga mag-aaral sa pamamagitan ng hands-on na pagkatuto at aplikasyon ng mga konsepto ng kimika sa mga sitwasyong hango sa totoong buhay. Nilalayon din ng proyektong ito na magsilbing kapaki-pakinabang na materyal para sa mga guro at mag-aaral upang mas mapadali ang pagkatuto ng kimika sa isang biswal, interaktibo, at makabuluhang paraan.

\vspace{1cm}

\noindent\textbf{Mga susing salita:} pisikal at kemikal na pagbabago, game-based learning, pang-edukasyon na laro
